% ---------------------------------------------------------------- packages
\usepackage[T1]{fontenc}
\usepackage{lmodern}
\usepackage{microtype}
\usepackage[margin=3.1cm]{geometry}

\usepackage{amsmath,amssymb,amsthm}
\usepackage{mathtools}
\usepackage{booktabs}
\usepackage{array}
\usepackage{caption}
\usepackage[ruled,vlined,linesnumbered]{algorithm2e}
\usepackage{tikz}
\usetikzlibrary{arrows.meta,positioning,calc,decorations.pathreplacing}
\usepackage{listings}
\usepackage[hidelinks]{hyperref}

\usepackage[backend=biber,style=alphabetic,maxbibnames=6]{biblatex}

% ---------------------------------------------------------------- theorems
\theoremstyle{plain}
\newtheorem{theorem}{Theorem}[section]

% ---------------------------------------------------------------- notation
\newcommand{\F}{\mathbb{F}_2}                  % the two-element field
\newcommand{\W}[1]{\mathsf{W}_{#1}}            % wiring by truth-table index
\newcommand{\wir}{w}                           % a generic wiring
\newcommand{\cxor}{\mathrm{cxor}}
\newcommand{\tape}{x}
\newcommand{\tab}{s}
\newcommand{\out}{r}
\newcommand{\addrset}{A}
\newcommand{\addr}[1]{a_{#1}}                  % the address in force at step #1
\newcommand{\cell}[1]{\tab[#1]}                % one bit of storage
\newcommand{\code}[1]{\texttt{#1}}
\newcommand{\bstr}[1]{\text{\texttt{#1}}}

% ---------------------------------------------------------------- listings
\lstset{basicstyle=\ttfamily\small, keywordstyle=\bfseries,
        commentstyle=\itshape, columns=fullflexible, frame=leftline,
        framesep=8pt, xleftmargin=10pt, showstringspaces=false, language=Python}
